\documentclass[11pt]{article}
\usepackage[a4paper,margin=25mm]{geometry}
\usepackage{amsmath,amssymb,amsthm,mathtools,mathrsfs,hyperref}
\newcommand{\Q}{\mathbb Q}
\newcommand{\Hd}{\operatorname{Hdg}}
\newcommand{\cl}{\operatorname{cl}}
\newcommand{\im}{\operatorname{im}}
\newcommand{\tr}{\operatorname{tr}}
\newcommand{\CH}{\operatorname{CH}}
\newcommand{\id}{\operatorname{id}}
\newcommand{\Span}{\operatorname{span}}
\newtheorem{theorem}{Theorem}
\newtheorem{proposition}[theorem]{Proposition}
\newtheorem{lemma}[theorem]{Lemma}
\title{GST universal source-existence attempt:\\
diagonal contraction and the one-sided source test}
\author{GST Hodge construction program}
\date{21 October 2026}
\begin{document}
\maketitle

\noindent\textbf{Proof verdict.}
The requested theorem is universal existence of actual cycles for the
canonical cycle-class map on every original smooth projective complex
variety. This note does \emph{not} prove that theorem.
It tests the new proposed forcing step from native diagonal,
duality, positive residual index, world closure, and formal alteration
trace identities to a one-sided native source. The explicit algebra
below disproves that inference from those specified abstract properties.
It is not a variety, not a canonical Chow theory, and not a counterexample
to the rational Hodge conjecture. It does not assert that all geometric
hypotheses of all GST theorems hold in the model.

\section{The exact theorem we tried to derive}
On the primitive middle rational Hodge space \(V_X\), use the established
positive form
\[
 Q_X(u,v)=(-1)^p\int_X uv,\qquad \dim X=2p.
\]
For a known actual native frame, let \(r\) be the prescribed orthogonal
residual. The previous blowup-world construction gives an actual
rational algebraic perfect-kernel class \(\mathscr Q\) and topological
pure-middle packets \(A_r,A_{Jr}\) satisfying
\[
 T_{\mathscr Q}A_r=A_r,\qquad
 -\chi_{Y\times Y}(A_r\boxtimes A_{Jr},\mathscr Q)=Q_X(r,r)>0.
 \tag{1}
\]
The desired new forcing assertion is
\[
 r\ne0\quad\Longrightarrow\quad
 \exists E\in K_0^{\rm alg}(Y)_{\Q}:
             \chi_Y(A_r,E)\ne0.                              \tag{U}
\]
The existing native Chow return converts such an \(E\) to an actual
cycle on the same original \(X\), with nonzero residual period.
Repeated rank increase would then terminate in finite dimension.
For a source-existence proof, (U) itself must be proved:
neither a factorization hypothesis nor a positive number is its proof.

The earlier derivation gives a correct constructor if an actual
one-sided factorization is supplied. We do not rename that conditional
constructor as an unconditional theorem. We now ask whether diagonal
contraction, duality and positivity force its missing input.

\section{An exact contraction theorem}
\begin{theorem}[Invariant tensor contraction preserves the source sector]
Let a rational vector space \(W\) have a linear involution \(g\) and
a nondegenerate symmetric bilinear form \(B\), with
\(B(gu,gv)=B(u,v)\). For \(n\geq0\), put
\[
 N_n=\{t\in W^{\otimes n}:g^{\otimes n}t=t\},\qquad N_0=\Q.
\]
Every finite expression formed from elements of the \(N_n\), rational
linear combinations, tensor products, permutations and contractions
by \(B\) has its \(m\)-leg output in \(N_m\).
This remains true when equivariant linear maps of arbitrary finite
arity are admitted as world-birth operations. In particular, if
\(gr=-r\), then every one-leg output \(z\) satisfies \(B(r,z)=0\).
\end{theorem}
\begin{proof}
Tensor products and permutations commute with the diagonal involution.
Invariance of \(B\) says precisely that contraction intertwines the
involutions on its input and output. An equivariant map sends fixed
vectors to fixed vectors. Induction on an expression proves the first
two assertions, without any bound on its finite arity or length.
For \(z\in N_1\),
\[
 B(r,z)=B(gr,gz)=B(-r,z)=-B(r,z),
\]
so \(B(r,z)=0\) over \(\Q\).
\end{proof}

The theorem allows \emph{all} invariant tensors, not just tensors
decomposable into known one-leg sources. It is therefore stronger than
a test of one selected finite frame. Numerical coefficients depending
on a prescribed target do not change the output space. Inserting a
nonnative target vector into a contraction, however, is not an
operation on native inputs; native preservation cannot be applied to it.

\section{A nonzero unital polarized graded test model}
For any \(a,b\in\Q_{>0}\), define the commutative graded algebra
\[
 H=\Q[h,s,r]/(h^5,hs,hr,sr,s^2-ah^4,r^2-bh^4),
 \quad \deg h=1,\quad\deg s=\deg r=2.                         \tag{2}
\]
Its basis and trace are
\[
 (1,h,h^2,h^3,h^4,s,r),\qquad
 \int h^4=1,\quad \int 1=\int h=\int h^2=\int h^3=\int s=\int r=0.
 \tag{3}
\]
Here degrees are codimensions. Give degree \(k\) the formal Tate
Hodge type \((k,k)\). Define \(B(u,v)=\int uv\).

\begin{lemma}
Equations (2)--(3) define a seven-dimensional graded Frobenius
algebra with Hard Lefschetz for \(h\) and positive primitive-middle
Hodge--Riemann form.
\end{lemma}
\begin{proof}
To see directly that the proposed basis has no hidden relation, define
multiplication on its seven-dimensional span by
\(h^ih^j=h^{i+j}\) for \(i+j\leq4\), zero otherwise,
\(1s=s,1r=r,s^2=ah^4,r^2=bh^4\), and zero for every other
product involving \(s\) or \(r\).
Associativity for three \(h\)-powers is truncated polynomial
associativity. With one primitive factor and two positive \(h\)-powers
both associations vanish. With two primitive factors and a positive
\(h\)-power, the only potentially nonzero inner product is a multiple
of \(h^4\), whose further product vanishes. With three primitive
factors both associations vanish. The unit cases are immediate.
This algebra satisfies the displayed relations, and those relations
reduce every monomial to the seven displayed elements. Thus (2) is
exactly this algebra.

On the \(h\)-chain, the matrix of \(B\) is the five-by-five
antidiagonal identity; its determinant is \(1\).
The remaining blocks are \(a,b\), so \(\det B=ab\ne0\).
The Lefschetz maps \(h^4:H^0\to H^4\) and
\(h^2:H^1\to H^3\) are isomorphisms, and the middle map is the
identity. The primitive middle space is
\[
 P^2=\ker(h:H^2\to H^3)=\Q s\oplus\Q r,\qquad
 B|_{P^2}=\begin{pmatrix}a&0\\0&b\end{pmatrix}>0.
\]
The primitive degree-zero form is \(\int h^4=1>0\);
the other primitive spaces in lower degree vanish.
These verify the even Tate Lefschetz and primitive sign conventions
of formal complex dimension four.
\end{proof}

Let \(g(h)=h,g(s)=s,g(r)=-r\). It is a trace-preserving graded
algebra automorphism. Define admissible tensors by the \(N_n\) of
Theorem 1, using this \(g\). The full one-leg sector is
\[
 A=N_1=\Q\{1,h,h^2,h^3,h^4,s\},\qquad
 A\cap P^2=\Q s.                                             \tag{4}
\]
It includes unit, ample class, top class and a nonzero primitive
class. Its restricted pairing is itself nondegenerate.
All formal rational Hodge classes belong to \(H\), including \(r\);
equation (4) still omits \(r\).

\section{The diagonal and positive residual survive}
The coevaluation tensor for \(B\) is
\[
 \Delta_B=\sum_{i=0}^4 h^i\otimes h^{4-i}
                  +\frac{s\otimes s}{a}+\frac{r\otimes r}{b}
                  \in N_2.                                  \tag{5}
\]
For \(C=\sum_i u_i\otimes v_i\), set
\[
 T_C(x)=\sum_i B(x,u_i)v_i.
\]
Then \(T_{\Delta_B}=\id_H\). The residual kernel
\[
 K=\frac{r\otimes r}{b}\in N_2
\]
satisfies
\[
 T_K^2=T_K,\quad T_K(r)=r,\quad T_K(A)=0,\quad
 \tr T_K=1,\quad
 (B\otimes B)(r\otimes r,K)=b>0.                              \tag{6}
\]
Thus an admissible projector onto a rational Tate residual exists
even though every admissible one-leg output annihilates that residual.

To check the exact sign of the established pure-middle Euler law,
take the \(p=2,M=5,J=\id\) convention on \(P^2\). Use
\(\chi=B\), the signed kernel \(\mathscr K=-K\), and
\[
 T^{\rm E}_C=-T_C,\qquad
 C\star_{\rm E}D=-C\star_BD .
\]
Here \(T_{C\star_BD}=T_C T_D\); the convention determines the
convolution. Then
\[
 \mathscr K\star_{\rm E}\mathscr K=\mathscr K,\quad
 T^{\rm E}_{\mathscr K}r=r,\quad
 -\chi_{\rm prod}(r\boxtimes r,\mathscr K)=b>0,\quad
 \chi(r,z)=0\quad(z\in A).                                  \tag{7}
\]
This matches the signs in (1). The model retains the positive doubled
test; it invalidates extraction of a one-leg source from that test.

There are six invariant and one anti-invariant basis vectors.
Consequently
\[
 \dim N_2=6^2+1^2=37,\qquad \dim(A\otimes A)=36.
\]
In fact \(N_2=(A\otimes A)\oplus\Q(r\otimes r)\).
The residual kernel lies exactly in the additional summand.
No one-sided factorization through \(A\) can reproduce its pairing
on \(r\otimes r\): each such expression is zero there, while (6)
is \(b\).

\section{Testing the other proposed forcing operations}
\paragraph{Native lifting.}
Every invariant two-leg kernel acts equivariantly and hence preserves
\(A\). Take the source carrier to be \(A\) and its realization the
inclusion \(c:A\hookrightarrow H\). Its native lift is simply the
restriction \(T|_A\); thus \(c(T|_A)=Tc\).
The class-to-defect block vanishes exactly as required by the GST
native-lifting theorem. For \(K\), this native lift is zero while its
ambient action on \(r\) is the identity.

Choose the \(B\)-orthogonal range retraction. Then the supported
extension of \(\id_A\) is the projection
\(P_A=\id_H-T_K\), rather than \(\id_H\). Both extensions agree on
\(A\). Equality on native inputs cannot identify them on \(r\).
This is exactly the extension freedom in
\texttt{supportedAmbientOperator\_id} and
\texttt{ambientNatural\_core\_decomposition}.

\paragraph{World closure and semantic reflection.}
Take roots in the invariant sectors and every birth map equivariant.
Theorem 1 proves invariant-sector stability at every finite stage,
including every finite arity. Hence semantic reflection applied to
the stable predicate \emph{lies in the invariant sector} preserves it.
It does not supply a world-generation premise for a noninvariant
one-leg vector. This does not negate any GST theorem which separately
assumes an actual surjective geometric assembly.

\paragraph{Degree and complete-alteration equations.}
For every positive integer \(e\), set \(H'=H^{\oplus e}\), with
diagonal involution and orthogonal direct-sum form. Define
\[
 f^*x=(x,\ldots,x),\qquad f_*(x_1,\ldots,x_e)=\sum_i x_i.
\]
Then
\[
 f_*f^*=e\,\id,\qquad B'(f^*x,f^*y)=eB(x,y),\qquad
 B'(f^*r,z')=0\quad(z'\in(H')^g).                            \tag{8}
\]
The last equality follows componentwise. In particular
\(B'(f^*r,f^*r)=eb>0\).
These are only the \emph{formal linear trace identities} of the
alteration argument; no actual morphism or geometric alteration is
being modeled. They show that the identities alone do not supply
a new source. Any geometry-specific forcing theorem needs more
than (8).

\section{An actual geometric check already inside GST}
There is also a genuine geometric test of the stronger idea
that an algebraic perfect diagonal automatically has finite separated
algebraic factors. Let \(C\) be a smooth projective curve of genus
\(g\geq1\), and on \(C\times C\) write
\[
 A=[\{c\}\times C],\quad B=[C\times\{c\}],\quad
 D=[\Delta_C]-A-B.
\]
The repository's existing all-genus curve intersection input gives
\[
 A^2=B^2=0,\ AB=1,\quad
 \Delta_C A=\Delta_C B=1,\quad \Delta_C^2=2-2g.
\]
Direct expansion yields
\[
 D A=D B=0,\qquad D^2=D\Delta_C=-2g\ne0.                     \tag{9}
\]
For any \(E,F\in K_0^{\rm alg}(C)_\Q\),
\[
 \operatorname{ch}_1(E\boxtimes F)=
 \operatorname{rk}(F)\operatorname{pr}_1^*\operatorname{ch}_1(E)
 +\operatorname{rk}(E)\operatorname{pr}_2^*\operatorname{ch}_1(F).
\]
In rational Betti cohomology this lies in \(\Q A+\Q B\), since
\(H^2(C,\Q)\) is generated by a point.
The sequence
\(0\to\mathcal O(-\Delta_C)\to\mathcal O\to\mathcal O_{\Delta_C}\to0\)
gives \(\operatorname{ch}_1(\mathcal O_{\Delta_C})=[\Delta_C]\).
If \([\mathcal O_{\Delta_C}]\) were a finite rational sum of external
\(K\)-products, its degree-one character would lie in
\(\Q A+\Q B\). Pairing with \(D\) contradicts (9).

This example excludes whole-diagonal separated \(K\)-factorization.
It does \emph{not} exclude the weaker factorization only on the
middle Hodge testing block used in the previous note. In particular
it is not a Hodge counterexample: \(D\) itself is an actual divisor.
The point is to distinguish two mathematically different constructors.


\section{Atomic scheme-level surgery: a genuine non-split source}
\noindent\textbf{Additional exact geometric test (21 October 2026).}
The invariant-sector model above exposes an algebraic logical mistake:
a nonzero doubled diagonal test does not generate a one-legged source.
We can strengthen that diagnosis with an actual smooth projective
complex variety and then exhibit the missing geometric input
\emph{as an independently constructed perfect-complex source}.
This test concerns a definite class of GST recipes; it neither assumes
nor proves universal Hodge algebraicity.

Let \(C\) be a smooth connected projective curve of genus
\(g\geq1\), and put \(X=C\times C\).
In \(\CH^1(X)_\Q\) take the actual fiber divisors
\(A=\{c\}\times C\), \(B=C\times\{c\}\), the ample
product polarization \(H=A+B\), and the internal graph
\[
  \delta=[\Delta_C],\qquad D=\delta-A-B.
\]
The normal bundle \(N_{\Delta_C/X}=T_C\) gives
\(\delta^2=\deg T_C=2-2g\), and intersection with
either fiber gives one. Consequently
\[
 A^2=B^2=0,\quad AB=1,\quad
 \delta A=\delta B=1,\quad
 H^2=2,\quad HD=0,\quad D^2=-2g.                  \tag{10}
\]
These are intersections of \emph{real Chow classes} on \(X\),
not abstract Frobenius-algebra coordinates. In particular
\(D\) is a nonzero algebraic primitive rational Hodge class,
with positive signed primitive norm \(2g\).

\begin{theorem}[Polarization--Chern--self-diagonal closure cannot birth \(D\)]
\label{th:tautological-no-go}
Consider the class of recipes starting with
\(\mathcal O_X,\mathcal O_X(H)\), and \(T_X\), and closed under finite
direct sums/differences in rational algebraic \(K_0\), tensor
products, duals, Adams operations and exterior powers of
honest vector-bundle expressions. Also allow algebraic
self-correspondence kernels \((\Delta_X)_*F\) built from the
same such \(F\) and applied to objects already in the recipe.
Every resulting one-legged \(K\)-class \(E\) satisfies
\[
       \cl(\ch_1(E))\in\Q h,\qquad
       \int_X\cl(\ch_1(E))\smile d=0,                \tag{11}
       \quad h=\cl(H),\quad d=\cl(D).
\]
No recipe in this class produces a class with
\(\ch_1(E)=D\), despite \(D\) being algebraic.
This holds at every finite operator-word length.
\end{theorem}
\begin{proof}
The initial characters have degree-one pieces
\[
 \cl(\ch_1(\mathcal O_X))=0,\qquad
 \cl(\ch_1(\mathcal O_X(H)))=h,\qquad
 \cl(\ch_1(T_X))=(2-2g)h .
\]
For any vector bundles (and by additivity for virtual classes),
\[
 \ch_1(E\oplus F)=\ch_1(E)+\ch_1(F),\qquad
 \ch_1(E\otimes F)=\rk(E)\ch_1(F)+\rk(F)\ch_1(E),
\]
\[
 \ch_1(E^\vee)=-\ch_1(E),\qquad
 \ch_1(\psi^kE)=k\ch_1(E),\qquad
 c_1(\wedge^kE)=\binom{\rk(E)-1}{k-1}c_1(E)
\]
for genuine rank-\(\rk(E)\) bundles where the exterior power
is defined. Each operation preserves the one-dimensional
\(\Q h\) first-character span \emph{after canonical Betti realization}.
For arbitrary \(C\), the equation
\(c_1(T_C)=(2-2g)[c]\) need NOT hold in
\(\CH^1(C)_\Q\): a degree-zero Picard-class difference can
survive rational equivalence while vanishing in Betti cohomology.
Only the Betti/numerical identity is used in this theorem.
For \(\Delta_X:X\to X\times X\), the projection formula gives
\[
 T_{(\Delta_X)_*F}(G)
  =(p_2)_*\bigl(p_1^*G\otimes(\Delta_X)_*F\bigr)
  =F\otimes G,                                      \tag{12}
\]
so the allowed self-diagonal kernels preserve it as well.
The diagonal self-intersection complex involves
\(\lambda_{-1}T_X^\vee\), and introduces no extra generator
beyond the stated characteristic operations.
Equation (10) proves (11).
If \(\cl(\ch_1(E))=d\), pairing with \(d\) would give
\(0=\int_X d^2=-2g\), contrary to \(g>0\).
Thus equality \(\ch_1(E)=D\) in Chow is also impossible.
\end{proof}

\begin{theorem}[Exact independent native repair by an internal graph]
\label{th:graph-repair}
Let \(L_D=\mathcal O_X(D)\) be the \emph{genuine} algebraic
line bundle of the internal diagonal-minus-fibers divisor.
The rational virtual perfect-complex class
\[
       \boxed{E_D=\tfrac12([L_D]-[L_D^{-1}])
               \in K_0^{\rm alg}(X)_\Q}
\]
has complete rational Chern character
\[
                  \ch(E_D)=D\in\CH^1(X)_\Q,
          \qquad -\int_X \ch_1(E_D)D=2g>0.           \tag{13}
\]
It therefore provides the algebraic one-leg source that
the entire restricted characteristic/self-diagonal recipe cannot.
\end{theorem}
\begin{proof}
Since \(\dim X=2\), the Chern-character expansions stop at
degree two:
\[
 \ch(L_D)=1+D+\frac{D^2}{2},\qquad
 \ch(L_D^{-1})=1-D+\frac{D^2}{2}.
\]
Their half-difference is precisely \(D\), with both rank and
degree-two characters zero. It is a rational combination
of actual line bundles, not a choice of target cycle.
Its numerical period is (10). This also correctly distinguishes
the diagonal \(\Delta_X\subset X\times X\), acting as the
identity/tensor kernel in (12), from the \emph{internal}
graph \(\Delta_C\subset C\times C=X\) supplying new geometry.
\end{proof}

At \(g=0\) the numerical graph norm is zero and no separating
period is obtained; for every \(g>0\) it is \(2g\).
The exact uniform intersection identities and genus-zero boundary
are kernel-checked in
\texttt{GSTPrimitiveProductTowerStabilization.lean}, beginning
with \texttt{curvePair\_tautological\_graph\_period\_zero} and
\texttt{curvePair\_internalGraph\_native\_period\_nonzero}.
The actual scheme-level interpretation relies on the displayed
curve normal-bundle and projective intersection identities.
A full arbitrary-\(X\) source-generation assertion is not
deduced from this family: it proves the need for genuinely
\emph{new geometry}, and provides an explicit successful non-split
test while eliminating an entire otherwise endlessly recycled
operator class.

\section{Source pins and executed checks}
All repository sources below were read at
\texttt{ae27ee80ae6c3448bd856a6b75bafd5f73543192}:
\begin{itemize}
\item \texttt{GSTClassicalHodgeNativeOperatorCohomologyRealization.lean}:\\
supported identity and ambient-extension decomposition.
\item \texttt{GSTClassicalHodgeKernelStableOperatorAlgebra.lean}:\\
native lift iff the class-to-defect block vanishes.
\item \texttt{GSTGraphV2OmniversalCore.lean}:\\
closure iff every seed-containing stable semantics holds.
\item \texttt{GSTGraphV2SelfExpandingCosmology.lean}:\\
generated-world theorem with actual assembly premises.
\item \texttt{GSTPrimitiveProductTowerStabilization.lean}:\\
the exact all-genus signed curve Gram.
\item The two previous derivations of native residual Euler kernels
and one-sided native factor extraction.
\end{itemize}
The exact Python script uses rational arithmetic at \(a=3,b=5\).
It checks all \(7^3=343\) basis associativity triples, all
\(7^2=49\) multiplication/automorphism pairs, nondegenerate pairing
with determinant \(15\), all three even Hard Lefschetz ranks,
primitive Gram \(\operatorname{diag}(3,5)\), projector identities,
invariant tensor dimensions and the signed doubled index \(-5\).
Trilinearity reduces associativity at this specialization to those
343 cases. Uniformity for all positive \(a,b\) is proved above,
not inferred from this specialization.

A separate Wolfram symbolic evaluation checks the primitive
contraction, inverse Gram, convolution sign, positive doubled index
\(b\), and (9) for symbolic \(a,b,g\). One preliminary evaluation
had a vector-flattening shape error; the corrected recorded evaluation
has no such messages. The computations verify the explicit models
and identities; they do not establish any canonical-cycle comparison
or universal geometric source existence.
All mathematical review here is self-review, not independent review.
No new Lean proof or compilation is claimed.

\medskip
\noindent\textbf{Remaining exact construction.}
A proof of (U) still needs a genuine algebraic object whose one-sided
residual period is nonzero. The contraction theorem proved here identifies
why diagonal duality, positive trace, native lifting and the formal
closure/degree identities cannot substitute for that object.
The original all-\(X\) construction goal remains unproved.
\end{document}

